\newpage
\subsubsection{Recursos disponibles}
El desarrollo del sistema de detección de insectos descortezadores
requiere la integración de un conjunto heterogéneo de tecnologías
que abarcan desde el entrenamiento de modelos de aprendizaje profundo
hasta el despliegue de infraestructura en la nube y el desarrollo de
aplicaciones móviles nativas. Las herramientas seleccionadas responden
a criterios de madurez tecnológica, compatibilidad entre componentes,
disponibilidad de licencias gratuitas para uso académico y adecuación
al enfoque offline-first del sistema.

En las siguientes subsecciones se describen las tecnologías empleadas
organizadas por capa del sistema: herramientas para el desarrollo del
modelo de clasificación, procesamiento de imágenes, desarrollo de la
aplicación móvil, desarrollo de la aplicación web e infraestructura
en la nube.